% ============================================================================
% SECTION 2: RELATED WORK
% ============================================================================
\section{Related Work}
\label{sec:related_work}

% ----------------------------------------------------------------------------
% 2.1 Adversarial Attacks
% ----------------------------------------------------------------------------
\subsection{Adversarial Attacks}
\label{subsec:attacks}

Adversarial attacks aim to find minimal perturbations that cause misclassification. The Fast Gradient Sign Method (FGSM)~\cite{goodfellow2015explaining} generates perturbations in a single step by following the sign of the loss gradient:
\begin{equation}
    x_{adv} = x + \eps \cdot \text{sign}(\nabla_x \mathcal{L}(f(x), y))
\end{equation}
where $\eps$ bounds the perturbation magnitude under the $\Linf$ norm.

Projected Gradient Descent (PGD)~\cite{madry2018towards} extends FGSM through iterative optimization:
\begin{equation}
    x^{t+1} = \Pi_{B_\eps(x)} \left( x^t + \alpha \cdot \text{sign}(\nabla_{x^t} \mathcal{L}(f(x^t), y)) \right)
\end{equation}
where $\Pi_{B_\eps(x)}$ projects onto the $\eps$-ball around the original input and $\alpha$ is the step size.

The Carlini \& Wagner (C\&W) attack~\cite{carlini2017towards} formulates adversarial example generation as an optimization problem, often achieving higher success rates but at increased computational cost. More recently, AutoAttack~\cite{croce2020reliable} combines multiple attack strategies (APGD-CE, APGD-DLR, FAB, Square Attack) to provide a reliable robustness evaluation, becoming the de facto standard for benchmarking.

% ----------------------------------------------------------------------------
% 2.2 Defense Mechanisms
% ----------------------------------------------------------------------------
\subsection{Defense Mechanisms}
\label{subsec:defenses}

Adversarial training (AT)~\cite{madry2018towards} remains the most effective defense, training models on adversarial examples generated during each batch:
\begin{equation}
    \min_\theta \mathbb{E}_{(x,y) \sim \mathcal{D}} \left[ \max_{\|\delta\| \leq \eps} \mathcal{L}(f_\theta(x + \delta), y) \right]
\end{equation}

TRADES~\cite{zhang2019theoretically} provides a theoretically-grounded trade-off between natural accuracy and robustness:
\begin{equation}
    \mathcal{L}_{\text{TRADES}} = \mathcal{L}_{CE}(f(x), y) + \beta \cdot \text{KL}(f(x) \| f(x_{adv}))
\end{equation}
where $\beta$ controls the trade-off parameter.

MART~\cite{wang2020improving} addresses the issue of misclassified examples by emphasizing them during training. Other approaches include input transformation~\cite{guo2018countering}, certified defenses via randomized smoothing~\cite{cohen2019certified}, and ensemble methods~\cite{tramer2018ensemble}. Recent advances in certified robustness include Adaptive Randomized Smoothing~\cite{ars2024neurips}, which improves certification radii through multi-step defenses.

% ----------------------------------------------------------------------------
% 2.3 CNN Robustness Studies
% ----------------------------------------------------------------------------
\subsection{CNN Robustness}
\label{subsec:cnn_robustness}

Extensive research has characterized the adversarial robustness of CNN architectures. Madry et al.~\cite{madry2018towards} demonstrated that adversarial training with PGD attacks provides strong empirical robustness. The RobustBench leaderboard~\cite{croce2021robustbench} tracks state-of-the-art CNN robustness, with WideResNet variants achieving the highest robust accuracy on CIFAR-10.

Studies have shown that model capacity correlates with achievable robustness~\cite{madry2018towards}, with wider and deeper networks demonstrating improved robust accuracy. Architecture-specific findings include the importance of skip connections~\cite{he2016deep} and batch normalization placement~\cite{xie2020adversarial}.

% ----------------------------------------------------------------------------
% 2.4 ViT Robustness Studies
% ----------------------------------------------------------------------------
\subsection{Vision Transformer Robustness}
\label{subsec:vit_robustness}

Vision Transformers~\cite{dosovitskiy2021image} process images as sequences of patches using self-attention. Initial studies suggested that ViTs exhibit some inherent robustness properties due to their global receptive field~\cite{bhojanapalli2021understanding}. However, subsequent work revealed that ViTs remain vulnerable to adversarial attacks, particularly gradient-based methods~\cite{mahmood2021robustness}.

Shao et al.~\cite{shao2022adversarial} demonstrated that adversarial training can improve ViT robustness, though challenges remain due to training instability. The relationship between patch size, attention mechanisms, and adversarial vulnerability has been explored~\cite{paul2022vision}. More recently, Jain et al.~\cite{jain2024towards} proposed Adaptive Attention Scaling Adversarial Training (AAS-AT) to address ViT-specific training challenges, achieving state-of-the-art robustness on CIFAR-10 and ImageNet-100. Wei et al.~\cite{wei2024enhancing} further investigated ViT vulnerabilities through gradient normalization and high-frequency adaptation. Despite these advances, comprehensive mechanistic comparisons between CNN and ViT robustness under unified evaluation frameworks remain limited.

% ----------------------------------------------------------------------------
% 2.5 Transfer Attacks
% ----------------------------------------------------------------------------
\subsection{Transfer Attacks}
\label{subsec:transfer}

Adversarial examples often transfer between models, enabling black-box attacks~\cite{papernot2016transferability}. The transferability depends on factors including model architecture, training data, and attack methods~\cite{demontis2019adversarial}.

Cross-architecture transfer between CNNs and ViTs has received increasing attention. Naseer et al.~\cite{naseer2022improving} explored transfer attacks from CNNs to ViTs, finding reduced transferability compared to within-architecture transfers. Recent work has proposed specialized methods for cross-architecture attacks: Token Gradient Regularization (TGR)~\cite{zhang2023tgr} crafts transferable adversarial samples by exploiting ViT's attention components, while Momentum Integrated Gradients (MIG)~\cite{ma2023transferable} achieves improved transferability across both ViTs and CNNs. Wang et al.~\cite{wang2023understanding} systematically analyze the factors affecting cross-architecture transferability, finding that architectural inductive biases---CNNs preferring high-frequency information while ViTs favor low-frequency patterns---significantly impact transfer success. Understanding these transfer patterns is crucial for developing robust ensemble systems and evaluating real-world attack scenarios.

% ----------------------------------------------------------------------------
% 2.6 Gap and Our Position
% ----------------------------------------------------------------------------
\subsection{Research Gap}
\label{subsec:gap}

While individual studies have examined CNN and ViT robustness separately, existing work primarily focuses on \textit{comparing} robustness scores rather than \textit{explaining} why architectures behave differently. Existing comparative studies typically report robustness metrics without mechanistic analysis, use inconsistent perturbation budgets or attack configurations across architectures, compare models with significantly different parameter counts, and focus on single aspects such as robustness or transferability in isolation.

Our work addresses these limitations by providing mechanistic explanations for observed differences. We investigate the gradient characteristics underlying transfer asymmetry, analyze layer-wise feature degradation in ViTs under adversarial perturbations, and connect these findings to design principles for robust systems. We employ the AutoAttack benchmark for reliable evaluation and demonstrate that gradient sparsity and alignment directly explain the observed behavioral differences.
